\documentclass[11pt,a4paper]{article}

% -------------------- Packages --------------------
\usepackage[
top=2cm,
bottom=2cm,
left=2cm,
right=2cm
]{geometry}

\usepackage{titlesec}
\usepackage{enumitem}
\usepackage{hyperref}
\usepackage{setspace}
\usepackage{graphicx}
\usepackage{amsmath}
\usepackage{amssymb}

% -------------------- Persian --------------------
\usepackage{xepersian}

\settextfont{XB Niloofar}
\setlatintextfont{TeX Gyre Termes}

\setstretch{1.1}

% -------------------- Lists --------------------
\setlist[itemize]{
	noitemsep,
	topsep=2pt
}

\setlist[enumerate]{
	noitemsep,
	topsep=2pt
}

% -------------------- Section formatting --------------------
\titleformat{\section}
{\large\bfseries}
{\thesection}
{0.7em}
{}

\titleformat{\subsection}
{\normalsize\bfseries}
{\thesubsection}
{0.7em}
{}

\titleformat{\subsubsection}
{\normalsize\bfseries}
{\thesubsubsection}
{0.7em}
{}

% -------------------- Hyperref --------------------
\hypersetup{
	colorlinks=true,
	linkcolor=blue,
	urlcolor=blue,
	citecolor=blue
}

% ==================================================
%                   DOCUMENT
% ==================================================

\begin{document}
	
	\begin{center}
		{\LARGE \textbf{پردازش تصویر دیجیتال}}\\[0.4cm]
		{\Large \textbf{جلسه دوم}}\\[0.3cm]
		{\large از «پیکسل» تا «بهبود تصویر»}\\[0.5cm]
		\textit{چگونه مقدار پیکسل‌ها را هدفمند تغییر دهیم تا تصویر برای هدف ما مناسب‌تر شود؟}
	\end{center}
	
	\vspace{0.3cm}
	
	% ==================================================
	\section{اهداف جلسه}
	
	در پایان این جلسه انتظار می‌رود دانشجو بتواند:
	
	\begin{itemize}
		\item \textbf{همسایگی، مجاورت و اتصال} پیکسل‌ها را تعریف کند و بین مجاورت ۴، ۸ و \lr{m} تمایز قائل شود.
		\item مفاهیم \textbf{مسیر، مؤلفه‌ی همبند، ناحیه و مرز} را توضیح دهد.
		\item \textbf{فاصله‌ی اقلیدسی، $D_4$ و $D_8$} را بین دو پیکسل محاسبه کند.
		\item \textbf{حوزه‌ی مکان} و رابطه‌ی $g(x,y)=T[f(x,y)]$ را توضیح دهد و عملیات نقطه‌ای را از عملیات همسایگی جدا کند.
		\item تبدیل‌های \textbf{منفی، لگاریتمی و توانی (گاما)} را به‌کار ببرد و اثر هر یک را تفسیر کند.
		\item \textbf{کشش کنتراست، آستانه‌گیری، برش سطح شدت و برش بیت‌پلین} را توضیح دهد.
		\item \textbf{هیستوگرام} یک تصویر را محاسبه کند و از شکل آن، وضعیت روشنایی و کنتراست را تشخیص دهد.
		\item \textbf{هم‌سازی هیستوگرام} را روی یک تصویر کوچک به‌صورت دستی اجرا کند.
		\item محدودیت‌های تبدیل‌های نقطه‌ای را بیان کند و توضیح دهد چرا به فیلترهای مکانی نیاز داریم.
	\end{itemize}
	
	\vspace{0.3cm}
	
	\noindent
	\textbf{ایده‌ی مرکزی جلسه:}
	
	\begin{quote}
		در جلسه‌ی قبل دیدیم که تصویر یک آرایه‌ی عددی است. اکنون می‌پرسیم:
		
		\begin{center}
			\textbf{«کدام تبدیل روی این اعداد، تصویر را برای هدف ما بهتر می‌کند؟»}
		\end{center}
		
		پاسخ به این سؤال همیشه به \textbf{کاربرد} بستگی دارد؛ یعنی «بهتر» بدون مشخص بودن هدف معنا ندارد.
	\end{quote}
	
	
	% ==================================================
	\section{مرور جلسه‌ی قبل و طرح مسئله}
	
	در جلسه‌ی اول دیدیم:
	
	\begin{itemize}
		\item تصویر دیجیتال تابع $f(x,y)$ با مختصات و مقدار گسسته است.
		\item با نمونه‌برداری ($M\times N$) و کوانتیزاسیون ($L=2^k$) ساخته می‌شود.
		\item مدل ساده‌ی تشکیل آن $f=i\,r$ است.
		\item بهبود تصویر فرایندی \textbf{ذهنی و وابسته به کاربرد} است.
	\end{itemize}
	
	حالا فرض کنید یک تصویر پرتو ایکس بیش از حد تیره است، یا تصویر یک صحنه کنتراست کمی دارد.
	چه کنیم؟
	
	برای پاسخ، دو ابزار پایه لازم داریم:
	
	\begin{enumerate}
		\item \textbf{زبانی برای صحبت درباره‌ی ارتباط پیکسل‌ها با یکدیگر} (همسایگی، مجاورت، فاصله)؛
		\item \textbf{تبدیل‌هایی که مقدار شدت را تغییر می‌دهند} (تبدیل‌های شدتی و هیستوگرام).
	\end{enumerate}
	
	این دو بخش اصلی جلسه‌ی امروز هستند.
	
	
	% ==================================================
	\section{روابط بنیادی بین پیکسل‌ها}
	
	یک تصویر مجموعه‌ای از پیکسل‌های جدا از هم نیست؛ پیکسل‌ها \textbf{در مکان کنار هم} قرار دارند.
	بسیاری از الگوریتم‌ها (لبه‌یابی، قطعه‌بندی، فیلتر کردن) روی همین ارتباط مکانی بنا شده‌اند.
	
	\subsection{همسایگی}
	
	برای پیکسل $p$ با مختصات $(x,y)$:
	
	\begin{itemize}
		\item \textbf{همسایه‌های چهارتایی} $N_4(p)$ (بالا، پایین، چپ، راست):
		\[
		(x+1,y),\quad(x-1,y),\quad(x,y+1),\quad(x,y-1).
		\]
		\item \textbf{همسایه‌های قطری} $N_D(p)$:
		\[
		(x+1,y+1),\quad(x+1,y-1),\quad(x-1,y+1),\quad(x-1,y-1).
		\]
		\item \textbf{همسایه‌های هشتتایی}:
		\[
		N_8(p)=N_4(p)\cup N_D(p).
		\]
	\end{itemize}
	
	اگر $p$ در لبه‌ی تصویر باشد، بعضی همسایه‌ها خارج از تصویر قرار می‌گیرند و باید جداگانه
	درباره‌ی آنها تصمیم بگیریم (حذف، تکرار مرز، صفر‌گذاری و ...).
	
	\subsection{مجاورت و اتصال}
	
	فرض کنید $V$ مجموعه‌ی مقدارهای شدتی باشد که برای تعریف «اتصال» به‌کار می‌بریم.
	مثلاً در تصویر دودویی ممکن است $V=\{1\}$ باشد.
	
	دو پیکسل $p$ و $q$ با مقدارهای در $V$ را:
	
	\begin{itemize}
		\item \textbf{مجاور-۴} می‌نامیم اگر $q\in N_4(p)$؛
		\item \textbf{مجاور-۸} می‌نامیم اگر $q\in N_8(p)$؛
		\item \textbf{مجاور-\lr{m}} (مجاورت مخلوط) می‌نامیم اگر:
		\begin{itemize}
			\item $q\in N_4(p)$، یا
			\item $q\in N_D(p)$ و مجموعه‌ی $N_4(p)\cap N_4(q)$ هیچ پیکسلی با مقدار در $V$ نداشته باشد.
		\end{itemize}
	\end{itemize}
	
	\textbf{چرا مجاورت \lr{m}؟} مجاورت-۸ گاهی \textbf{ابهام} ایجاد می‌کند: ممکن است بین دو پیکسل
	بیش از یک مسیر وجود داشته باشد. مجاورت \lr{m} با حذف اتصال قطری «اضافی»، این ابهام را برطرف می‌کند.
	
	\subsection{مسیر، مؤلفه‌ی همبند، ناحیه و مرز}
	
	\begin{itemize}
		\item \textbf{مسیر} از $p$ به $q$: دنباله‌ای از پیکسل‌های متمایز که هر دو عضو متوالی آن مجاورند.
		\item \textbf{مؤلفه‌ی همبند:} مجموعه‌ی پیکسل‌هایی از یک زیرمجموعه $S$ که بین هر دو تای آنها یک مسیر درون $S$ وجود دارد.
		\item \textbf{ناحیه:} زیرمجموعه‌ای همبند از تصویر. دو ناحیه را \textbf{مجاور} می‌نامیم اگر اجتماع آنها همبند باشد.
		\item \textbf{مرز} یک ناحیه: پیکسل‌هایی از ناحیه که حداقل یک همسایه خارج از ناحیه دارند.
	\end{itemize}
	
	\begin{quote}
		\textbf{تعریف «همبند بودن» به نوع مجاورتی که انتخاب می‌کنیم وابسته است؛
			پس تعداد مؤلفه‌ها با مجاورت-۴ و مجاورت-۸ ممکن است فرق کند.}
	\end{quote}
	
	
	% ==================================================
	\section{مثال: مجاورت در یک تصویر کوچک}
	
	تصویر دودویی زیر را با $V=\{1\}$ در نظر بگیرید (سطر و ستون از صفر شماره‌گذاری شده‌اند):
	
	\[
	\begin{bmatrix}
		0 & 1 & 1\\
		0 & 1 & 0\\
		0 & 0 & 1
	\end{bmatrix}.
	\]
	
	پیکسل‌های با مقدار $1$ در مختصات زیر هستند:
	
	\[
	(0,1),\quad(0,2),\quad(1,1),\quad(2,2).
	\]
	
	\begin{itemize}
		\item پیکسل‌های $(0,2)$ و $(1,1)$ با مجاورت-۸ مجاورند (قطری)، اما از راه $(0,1)$ هم
		می‌توان از یکی به دیگری رفت. یعنی \textbf{دو مسیر} داریم؛ این همان ابهام است.
		\item با مجاورت \lr{m}، پیکسل‌های $(0,2)$ و $(1,1)$ مجاور نیستند؛ چون
		$N_4$ مشترک آنها شامل $(0,1)$ با مقدار $1$ است.
		\item پیکسل‌های $(1,1)$ و $(2,2)$ با مجاورت \lr{m} مجاورند؛ زیرا
		$N_4$ مشترک آنها فقط شامل پیکسل‌هایی با مقدار $0$ است.
	\end{itemize}
	
	
	% ==================================================
	\section{معیارهای فاصله}
	
	برای پیکسل‌های $p=(x,y)$ و $q=(s,t)$ سه فاصله‌ی رایج داریم:
	
	\[
	D_e(p,q)=\sqrt{(x-s)^2+(y-t)^2}
	\qquad\text{(اقلیدسی)}
	\]
	
	\[
	D_4(p,q)=|x-s|+|y-t|
	\qquad\text{(بلوک شهری)}
	\]
	
	\[
	D_8(p,q)=\max\big(|x-s|,\,|y-t|\big)
	\qquad\text{(صفحه‌ی شطرنج)}
	\]
	
	هر تابع فاصله‌ی $D$ باید سه شرط را داشته باشد:
	
	\begin{enumerate}
		\item $D(p,q)\geq0$ و $D(p,q)=0$ اگر و فقط اگر $p=q$؛
		\item $D(p,q)=D(q,p)$؛
		\item $D(p,z)\leq D(p,q)+D(q,z)$ (نامساوی مثلث).
	\end{enumerate}
	
	شکل مجموعه‌ی نقاطی که فاصله‌شان از یک پیکسل حداکثر $r$ است:
	
	\begin{itemize}
		\item برای $D_e$: یک \textbf{دایره}؛
		\item برای $D_4$: یک \textbf{لوزی}؛
		\item برای $D_8$: یک \textbf{مربع}.
	\end{itemize}
	
	\textbf{مثال:} برای $p=(2,3)$ و $q=(6,6)$:
	
	\[
	D_e=\sqrt{4^2+3^2}=5,
	\qquad
	D_4=4+3=7,
	\qquad
	D_8=\max(4,3)=4.
	\]
	
	نکته: $D_4$ و $D_8$ تنها به مختصات دو پیکسل بستگی دارند؛ اما اگر بخواهیم فاصله را
	\textbf{روی یک مسیر مجاز} (مثلاً فقط از پیکسل‌های با مقدار $1$) اندازه بگیریم،
	فاصله به مقدار پیکسل‌های میانی هم وابسته می‌شود.
	
	
	% ==================================================
	\section{حوزه‌ی مکان و عملیات روی تصویر}
	
	بخش بزرگی از پردازش تصویر در \textbf{حوزه‌ی مکان} انجام می‌شود؛ یعنی مستقیماً روی خود پیکسل‌ها:
	
	\[
	g(x,y)=T\big[f(x,y)\big].
	\]
	
	که در آن $f$ تصویر ورودی، $g$ تصویر خروجی و $T$ یک عملگر است که روی همسایگی
	نقطه‌ی $(x,y)$ تعریف می‌شود.
	
	دو حالت اصلی:
	
	\begin{itemize}
		\item \textbf{عملیات نقطه‌ای:} همسایگی فقط یک پیکسل است؛ پس با $r=f(x,y)$ و $s=g(x,y)$:
		\[
		s=T(r).
		\]
		به این تبدیل‌ها \textbf{تبدیل‌های شدتی} می‌گوییم. این جلسه روی همین حالت متمرکز است.
		\item \textbf{عملیات همسایگی:} مقدار خروجی به پیکسل‌های اطراف هم وابسته است
		(فیلترهای مکانی، موضوع جلسه‌ی بعد).
	\end{itemize}
	
	یادآوری: در پردازش تصویر، اگر خلاف آن گفته نشود، عملیات بین دو تصویر
	(جمع، ضرب و ...) به‌صورت \textbf{درایه به درایه} انجام می‌شود، نه به شکل ضرب ماتریسی.
	
	
	% ==================================================
	\section{تبدیل‌های شدتی پایه}
	
	در این بخش فرض می‌کنیم شدت‌ها در بازه‌ی $[0,L-1]$ هستند.
	
	\subsection{منفی تصویر}
	
	\[
	s=L-1-r.
	\]
	
	سیاه به سفید و سفید به سیاه تبدیل می‌شود. برای نمایش جزئیات روشن یا سفید در یک زمینه‌ی تیره‌ی بزرگ مفید است.
	
	\subsection{تبدیل لگاریتمی}
	
	\[
	s=c\,\log(1+r),
	\qquad r\geq0.
	\]
	
	\begin{itemize}
		\item بازه‌ی باریکی از مقدارهای \textbf{تیره} را به بازه‌ی پهن‌تری از خروجی می‌نگارد (گسترش تیره‌ها).
		\item مقدارهای روشن را \textbf{فشرده} می‌کند.
		\item برای نمایش داده‌هایی با دامنه‌ی دینامیکی بسیار بزرگ، مثل طیف فوریه، کاربرد دارد.
	\end{itemize}
	
	برای این‌که بازه‌ی $[0,L-1]$ به خودش نگاشته شود:
	
	\[
	c=\frac{L-1}{\ln L}.
	\]
	
	\textbf{مثال:} برای $L=256$ داریم $c=\dfrac{255}{\ln256}\approx45.99$. در نتیجه:
	
	\[
	r=0\Rightarrow s=0,
	\qquad
	r=15\Rightarrow s\approx45.99\times\ln16\approx128,
	\qquad
	r=255\Rightarrow s=255.
	\]
	
	یعنی همه‌ی مقدارهای از $0$ تا $15$ (کمتر از ۶ درصد بازه) تقریباً نیمی از دامنه‌ی خروجی را اشغال می‌کنند.
	
	\subsection{تبدیل توانی (گاما)}
	
	\[
	s=c\,r^{\gamma}.
	\]
	
	در حالت نرمال‌شده، یعنی $r\in[0,1]$ و $c=1$، خروجی نیز در $[0,1]$ می‌ماند:
	
	\begin{itemize}
		\item $\gamma<1$: مقدارهای تیره گسترش می‌یابند، شبیه لگاریتم؛ تصویر \textbf{روشن‌تر} می‌شود.
		\item $\gamma>1$: مقدارهای روشن گسترش می‌یابند؛ تصویر \textbf{تیره‌تر} می‌شود.
		\item $\gamma=1$: تبدیل همانی.
	\end{itemize}
	
	برخلاف لگاریتم که یک منحنی ثابت است، با تغییر $\gamma$ خانواده‌ای از منحنی‌ها داریم.
	
	\textbf{کاربرد مهم: تصحیح گاما.} بسیاری از نمایشگرها رابطه‌ی ورودی--خروجی‌شان تقریباً از نوع
	توانی با $\gamma$ حدود ۱٫۸ تا ۲٫۵ است و تصویر را تیره‌تر از حالت مطلوب نشان می‌دهند.
	برای جبران، پیش از نمایش، تبدیلی با $\gamma$ برابر معکوس آن اعمال می‌کنیم.
	
	\textbf{مثال:} برای $L=256$ و $r=64$ (نرمال‌شده: $0.251$):
	
	\[
	\gamma=0.5:\quad s=255\times0.251^{0.5}\approx128,
	\qquad
	\gamma=2:\quad s=255\times0.251^{2}\approx16.
	\]
	
	یعنی همان پیکسل با $\gamma=0.5$ بسیار روشن‌تر و با $\gamma=2$ بسیار تیره‌تر می‌شود.
	
	
	% ==================================================
	\section{تبدیل‌های خطی تکه‌ای}
	
	گاهی به‌جای یک فرمول بسته، یک تابع تکه‌ای ساده با نقاط شکست مشخص تعریف می‌کنیم.
	مزیت آن \textbf{انعطاف زیاد} و عیب آن \textbf{نیاز به پارامترهای بیشتر} است.
	
	\subsection{کشش کنتراست}
	
	اگر شدت‌های تصویر فقط در بازه‌ی کوچکی از $[0,L-1]$ قرار گرفته باشند،
	کنتراست کم است. با تبدیل خطی می‌توان بازه را کشید:
	
	\[
	s=(L-1)\,\frac{r-r_{\min}}{r_{\max}-r_{\min}}.
	\]
	
	\textbf{مثال:} تصویری با $L=256$ که شدت‌هایش بین $90$ و $150$ است:
	
	\[
	r=90\Rightarrow s=0,
	\qquad
	r=120\Rightarrow s=255\times\frac{30}{60}\approx128,
	\qquad
	r=150\Rightarrow s=255.
	\]
	
	حالت کلی‌تر آن، تبدیلی با دو نقطه‌ی کنترل $(r_1,s_1)$ و $(r_2,s_2)$ است.
	
	\subsection{آستانه‌گیری}
	
	حالت حدی کشش کنتراست:
	
	\[
	s=
	\begin{cases}
		L-1 & r>T\\
		0 & r\leq T
	\end{cases}
	\]
	
	خروجی یک \textbf{تصویر دودویی} است. آستانه‌گیری ساده‌ترین شکل \textbf{قطعه‌بندی} است
	و در جلسه‌های بعد دوباره به آن برمی‌گردیم.
	
	\subsection{برش سطح شدت}
	
	هدف برجسته کردن یک بازه‌ی مشخص از شدت‌هاست، مثلاً $[A,B]$:
	
	\begin{itemize}
		\item روش اول: مقدارهای داخل بازه را روشن و بقیه را تیره می‌کنیم.
		\item روش دوم: مقدارهای داخل بازه را روشن می‌کنیم و بقیه را بدون تغییر نگه می‌داریم.
	\end{itemize}
	
	کاربرد: برجسته کردن یک نوع بافت یا ماده در تصویر پزشکی و صنعتی.
	
	\subsection{برش بیت‌پلین}
	
	در تصویر $8$ بیتی، هر پیکسل با $8$ بیت نمایش داده می‌شود. می‌توانیم هر بیت را جدا کنیم
	و $8$ «صفحه‌ی بیتی» بسازیم.
	
	\textbf{مثال:} عدد $200$ در مبنای دو:
	
	\[
	200=128+64+8=(11001000)_2.
	\]
	
	\begin{itemize}
		\item بیت‌های \textbf{پرارزش} (مثلاً سه بیت اول) بیشتر اطلاعات ساختاری و شکل تصویر را حمل می‌کنند.
		\item بیت‌های \textbf{کم‌ارزش} بیشتر جزئیات ریز و نویز مانند را دارند.
	\end{itemize}
	
	این موضوع هم برای \textbf{فشرده‌سازی} (نگه داشتن بیت‌های مهم‌تر) و هم برای
	شناخت اهمیت هر بیت مفید است.
	
	
	% ==================================================
	\section{هیستوگرام تصویر}
	
	\textbf{هیستوگرام} یک تصویر با شدت‌های $r_k$ (برای $k=0,1,\ldots,L-1$) برابر است با:
	
	\[
	h(r_k)=n_k,
	\]
	
	که در آن $n_k$ تعداد پیکسل‌هایی است که شدتشان $r_k$ است.
	
	هیستوگرام \textbf{نرمال‌شده}:
	
	\[
	p(r_k)=\frac{n_k}{MN},
	\qquad
	\sum_{k=0}^{L-1}p(r_k)=1.
	\]
	
	می‌توان $p(r_k)$ را \textbf{برآوردی از احتمال} وجود شدت $r_k$ در تصویر دانست.
	
	\textbf{شکل هیستوگرام چه می‌گوید؟}
	
	\begin{itemize}
		\item \textbf{تصویر تیره:} هیستوگرام نزدیک چپ (شدت‌های کم) متمرکز است.
		\item \textbf{تصویر روشن:} هیستوگرام نزدیک راست (شدت‌های زیاد) است.
		\item \textbf{کنتراست کم:} هیستوگرام باریک و در مرکز بازه است.
		\item \textbf{کنتراست زیاد و خوب:} هیستوگرام در بازه‌ی وسیعی پخش شده است.
	\end{itemize}
	
	نکته‌ی مهم این است که هیستوگرام \textbf{اطلاعات مکانی ندارد}: دو تصویر کاملاً متفاوت
	می‌توانند هیستوگرام یکسان داشته باشند. با این حال، بسیار ارزان محاسبه می‌شود
	و برای ارزیابی و بهبود تصویر بسیار مفید است.
	
	
	% ==================================================
	\section{هم‌سازی هیستوگرام}
	
	ایده: تبدیلی پیدا کنیم که هیستوگرام خروجی را \textbf{تا حد ممکن یکنواخت} کند؛
	یعنی همه‌ی سطوح شدت تقریباً به یک اندازه ظاهر شوند و کنتراست افزایش یابد.
	
	\subsection{حالت پیوسته}
	
	فرض کنید $r\in[0,L-1]$ یک متغیر پیوسته با چگالی $p_r(r)$ است. تبدیل زیر را در نظر بگیرید:
	
	\[
	s=T(r)=(L-1)\int_{0}^{r}p_r(w)\,dw.
	\]
	
	این تابع (که مضربی از \textbf{تابع توزیع تجمعی} است) یکنوا صعودی است
	و بازه‌ی $[0,L-1]$ را به همان بازه می‌نگارد. چگالی خروجی از رابطه‌ی
	
	\[
	p_s(s)=p_r(r)\left|\frac{dr}{ds}\right|
	\]
	
	به دست می‌آید. چون
	
	\[
	\frac{ds}{dr}=(L-1)\,p_r(r),
	\]
	
	پس:
	
	\[
	p_s(s)=p_r(r)\cdot\frac{1}{(L-1)\,p_r(r)}=\frac{1}{L-1}.
	\]
	
	یعنی خروجی دقیقاً \textbf{یکنواخت} است، صرف‌نظر از شکل $p_r$.
	
	\subsection{حالت گسسته}
	
	در تصویر دیجیتال:
	
	\[
	s_k=T(r_k)=(L-1)\sum_{j=0}^{k}p(r_j)
	=\frac{L-1}{MN}\sum_{j=0}^{k}n_j,
	\]
	
	و مقدار حاصل را به نزدیک‌ترین عدد صحیح گرد می‌کنیم.
	
	توجه کنید که در حالت گسسته هیستوگرام خروجی \textbf{دقیقاً مسطح نمی‌شود}؛
	فقط معمولاً پهن‌تر می‌شود و سطوح شدت از هم فاصله می‌گیرند.
	
	
	% ==================================================
	\section{یک مثال کامل: هم‌سازی هیستوگرام}
	
	تصویر $64\times64$ با $3$ بیت (یعنی $L=8$ و $MN=4096$) و هیستوگرام زیر را در نظر بگیرید:
	
	\[
	\begin{array}{c|cccccccc}
		r_k & 0 & 1 & 2 & 3 & 4 & 5 & 6 & 7\\
		\hline
		n_k & 790 & 1023 & 850 & 656 & 329 & 245 & 122 & 81\\
		\sum n_j & 790 & 1813 & 2663 & 3319 & 3648 & 3893 & 4015 & 4096\\
		7\sum n_j/4096 & 1.35 & 3.10 & 4.55 & 5.67 & 6.23 & 6.65 & 6.86 & 7.00\\
		s_k & 1 & 3 & 5 & 6 & 6 & 7 & 7 & 7
	\end{array}
	\]
	
	سطرهای جدول به ترتیب: هیستوگرام، جمع تجمعی، مقدار قبل از گرد کردن، و شدت جدید.
	
	هیستوگرام جدید با جمع کردن پیکسل‌هایی که به یک مقدار نگاشته شده‌اند:
	
	\[
	\begin{array}{c|ccccc}
		s & 1 & 3 & 5 & 6 & 7\\
		\hline
		\text{تعداد} & 790 & 1023 & 850 & 985 & 448
	\end{array}
	\]
	
	بررسی: $790+1023+850+985+448=4096$.
	
	مشاهدات:
	
	\begin{itemize}
		\item شدت‌های ۰، ۲ و ۴ دیگر ظاهر نمی‌شوند؛ هیستوگرام \textbf{مسطح نشده}، اما پخش‌تر شده است.
		\item سطوح ۶ و ۷ از ترکیب چند سطح قبلی ساخته شدند؛ \textbf{اطلاعات از دست رفته است}
		(تبدیل یکبه‌یک نیست).
		\item تصویر قبلی مقادیر کم و تیره زیادی داشت؛ اکنون بازه‌ی دینامیکی بزرگ‌تری دارد.
	\end{itemize}
	
	
	% ==================================================
	\section{هم‌سازی هیستوگرام: نکات و محدودیت‌ها}
	
	\begin{itemize}
		\item \textbf{خودکار است:} بر خلاف گاما، پارامتر ندارد؛ اما همین ممکن است خوب یا بد باشد.
		\item \textbf{سراسری است:} تبدیل برای کل تصویر یکسان است و ممکن است جزئیات محلی را خوب نکند.
		نسخه‌های \textbf{محلی} (مانند \lr{CLAHE}) هیستوگرام را در پنجره‌های کوچک محاسبه می‌کنند.
		\item \textbf{نویز را تقویت می‌کند:} اگر ناحیه‌ای از تصویر محتوای کمی دارد ولی نویز دارد،
		بعد از هم‌سازی نویز آشکارتر می‌شود.
		\item \textbf{همیشه نتیجه‌ی بهتر نمی‌دهد:} گاهی تصویر «طبیعی‌تر» با هیستوگرام غیر یکنواخت است.
		\item \textbf{هیستوگرام مطلوب دیگر:} اگر بخواهیم هیستوگرام خروجی شکل خاصی داشته باشد
		(نه یکنواخت)، از \textbf{تطبیق هیستوگرام} (\lr{Histogram Matching}) استفاده می‌کنیم.
	\end{itemize}
	
	\begin{quote}
		\textbf{هم‌سازی هیستوگرام یک ابزار است، نه یک پاسخ جادویی برای همه‌ی تصویرها.}
	\end{quote}
	
	
	% ==================================================
	\section{کدام روش، برای کدام مسئله؟}
	
	\begin{itemize}
		\item تصویر بسیار تیره، جزئیات در تیره‌ها: \textbf{گاما با $\gamma<1$} یا \textbf{لگاریتم}.
		\item تصویر بیش از حد روشن: \textbf{گاما با $\gamma>1$}.
		\item جزئیات روشن در پس‌زمینه‌ی تیره: \textbf{منفی تصویر}.
		\item داده با دامنه‌ی دینامیکی بسیار وسیع (مثل طیف فوریه): \textbf{لگاریتم}.
		\item شدت‌ها در بازه‌ی باریک: \textbf{کشش کنتراست} یا \textbf{هم‌سازی هیستوگرام}.
		\item جداسازی شیء از پس‌زمینه: \textbf{آستانه‌گیری}.
		\item برجسته کردن یک بازه‌ی شدتی: \textbf{برش سطح شدت}.
		\item نمایش درست روی نمایشگر: \textbf{تصحیح گاما}.
	\end{itemize}
	
	این‌ها \textbf{راهنمایی کلی} هستند؛ انتخاب نهایی معمولاً با آزمون و ارزیابی نتیجه انجام می‌شود،
	چون بهبود تصویر ذاتاً به کاربرد وابسته است.
	
	
	% ==================================================
	\section{یک محدودیت بنیادی: تبدیل‌های نقطه‌ای کافی نیستند}
	
	تمام تبدیل‌های امروز $s=T(r)$ برای هر پیکسل \textbf{مستقل} از همسایه‌ها عمل می‌کنند.
	در نتیجه:
	
	\begin{itemize}
		\item نمی‌توانند \textbf{نویز} را کم کنند؛ چون برای تشخیص نویز باید پیکسل را با اطرافش مقایسه کرد.
		\item نمی‌توانند \textbf{لبه‌ها} را تیز یا آشکار کنند؛ چون لبه یعنی تغییر شدت بین پیکسل‌های مجاور.
		\item دو پیکسل با شدت یکسان، همیشه خروجی یکسان می‌گیرند، حتی اگر یکی در ناحیه‌ی هموار و دیگری روی لبه باشد.
	\end{itemize}
	
	برای حل این مسائل باید از \textbf{اطلاعات همسایگی} استفاده کنیم؛ همان چیزی که در
	بخش روابط پیکسل‌ها تعریف کردیم. این موضوع جلسه‌ی بعد است.
	
	
	% ==================================================
	\section{جمع‌بندی مفهومی جلسه}
	
	\[
	\boxed{
		\text{تصویر}
		\rightarrow
		\text{روابط مکانی پیکسل‌ها}
		\rightarrow
		\text{تبدیل شدتی}
		\rightarrow
		\text{هیستوگرام}
		\rightarrow
		\text{تصویر بهبودیافته}
	}
	\]
	
	مهم‌ترین اصطلاحات جلسه:
	
	\begin{itemize}
		\item \textbf{همسایگی} $N_4$، $N_D$ و $N_8$:
		مجموعه‌ی پیکسل‌های اطراف یک پیکسل.
		
		\item \textbf{مجاورت} (۴، ۸ و \lr{m}):
		قواعد تعیین این‌که دو پیکسل به هم متصل‌اند یا نه.
		
		\item \textbf{مؤلفه‌ی همبند و ناحیه:}
		مجموعه‌ی پیکسل‌های به‌هم‌پیوسته.
		
		\item \textbf{فاصله‌ی $D_e$ ،$D_4$ و $D_8$:}
		سه معیار برای اندازه‌گیری دوری پیکسل‌ها.
		
		\item \textbf{عملیات نقطه‌ای:}
		$s=T(r)$ که فقط به همان پیکسل بستگی دارد.
		
		\item \textbf{تبدیل لگاریتمی و توانی:}
		گسترش تیره‌ها یا روشن‌ها با منحنی‌های غیرخطی.
		
		\item \textbf{کشش کنتراست و آستانه‌گیری:}
		تبدیل‌های خطی تکه‌ای.
		
		\item \textbf{هیستوگرام:}
		$h(r_k)=n_k$، شمارش تعداد پیکسل‌ها در هر سطح شدت.
		
		\item \textbf{هم‌سازی هیستوگرام:}
		$s_k=(L-1)\sum_{j\leq k}p(r_j)$.
	\end{itemize}
	
	
	% ==================================================
	\section{تمرین کلاسی ۱: همبندی و فاصله}
	
	تصویر دودویی زیر را با $V=\{1\}$ در نظر بگیرید:
	
	\[
	\begin{bmatrix}
		0 & 1 & 1 & 0 & 0\\
		1 & 1 & 0 & 0 & 1\\
		0 & 0 & 1 & 1 & 1\\
		0 & 1 & 0 & 0 & 0\\
		1 & 1 & 0 & 1 & 0
	\end{bmatrix}.
	\]
	
	\begin{enumerate}
		\item تعداد مؤلفه‌های همبند را با مجاورت-۴ و با مجاورت-۸ بشمارید. چرا فرق دارند؟
		\item آیا پیکسل‌های $(1,1)$ و $(2,2)$ با مجاورت \lr{m} مجاورند؟ دلیل بیاورید.
		\item برای $p=(1,1)$ و $q=(7,9)$ مقدارهای $D_e$ و $D_4$ و $D_8$ را حساب کنید.
		\item مرز مؤلفه‌ی شامل $(2,2)$ را (با مجاورت-۸) مشخص کنید.
	\end{enumerate}
	
	
	% ==================================================
	\section{تمرین کلاسی ۲: انتخاب تبدیل مناسب}
	
	برای هر وضعیت، یک تبدیل (یا ترکیبی از آنها) پیشنهاد دهید و دلیل خود را بنویسید:
	
	\begin{enumerate}
		\item عکس چهره که هنگام غروب با نور کم گرفته شده و بسیار تیره است؛
		\item تصویر ماموگرافی که توده‌ی روشن در زمینه‌ی تیره دارد؛
		\item نمایش طیف قدرت فوریه‌ی یک تصویر که چند مقدار بسیار بزرگ دارد؛
		\item تصویر مه‌آلود با شدت‌هایی در بازه‌ی $[100,160]$ از $[0,255]$؛
		\item تصویر یک متن اسکن‌شده که می‌خواهیم حروف و کاغذ را از هم جدا کنیم؛
		\item نمایش درست یک تصویر روی نمایشگری با $\gamma=2.2$.
	\end{enumerate}
	
	
	% ==================================================
	\section{تمرین ریاضی}
	
	\begin{enumerate}
		\item برای $L=256$ مقدار $c$ در تبدیل $s=c\ln(1+r)$ را به‌گونه‌ای بیابید که $r=255$ به $s=255$ برود.
		سپس $s$ را برای $r=0$، $r=7$ و $r=31$ حساب کنید.
		\item برای $L=256$ و $r=100$ مقدار خروجی تبدیل توانی نرمال‌شده را برای
		$\gamma=0.4$ و $\gamma=2.5$ حساب کنید و نتیجه را تفسیر کنید.
		\item تصویری $8$ بیتی دارای شدت‌هایی بین $60$ و $180$ است. تبدیل کشش کنتراست آن
		را بنویسید و $s$ را برای $r=60$، $r=100$ و $r=180$ بیابید.
		\item عدد $173$ را به مبنای دو ببرید. اگر فقط سه بیت پرارزش را نگه داریم و بقیه را صفر کنیم، مقدار جدید چه می‌شود؟
		\item تصویر $3$ بیتی‌ای با $16$ پیکسل دارای هیستوگرام زیر است:
		\[
		\begin{array}{c|cccccccc}
			r_k & 0 & 1 & 2 & 3 & 4 & 5 & 6 & 7\\
			\hline
			n_k & 2 & 4 & 3 & 1 & 0 & 2 & 3 & 1
		\end{array}
		\]
		جدول هم‌سازی هیستوگرام (مقدار $s_k$ و هیستوگرام جدید) را کامل کنید.
		\item نشان دهید که تبدیل هم‌سازی \textbf{یکنوا و غیرنزولی} است؛ یعنی $s_{k+1}\geq s_k$. (راهنمایی: $p(r_k)\geq0$.)
	\end{enumerate}
	
	
	% ==================================================
	\section{تمرین فکری}
	
	\begin{enumerate}
		\item آیا هم‌سازی هیستوگرام یا گاما می‌تواند \textbf{اطلاعات جدید} به تصویر اضافه کند؟
		یا فقط اطلاعات موجود را آشکارتر می‌کند؟ (به مثال ۳ بیتی جلسه دقت کنید: چند سطح شدت پس از هم‌سازی باقی ماند؟)
		\item چرا دو تصویر متفاوت ممکن است هیستوگرام یکسان داشته باشند؟ یک مثال کوچک بزنید.
		\item چرا هم‌سازی هیستوگرام روی یک تصویر که بیشتر آن پس‌زمینه‌ی یکدست و فقط بخش کوچکی شیء است،
		ممکن است نتیجه‌ی بدی بدهد؟
		\item در چه شرایطی \textbf{مجاورت-۸} ممکن است نتیجه‌ی غیرمنتظره بدهد و مجاورت \lr{m} آن را برطرف کند؟
		\item با توجه به اینکه $T$ برای هر پیکسل مستقل عمل می‌کند، چه کاری را
		هیچ تبدیل نقطه‌ای نمی‌تواند انجام دهد؟
	\end{enumerate}
	
	
	% ==================================================
	\section{تمرین عملی (در نرم‌افزار)}
	
	با \lr{Python} (کتابخانه‌هایی مانند \lr{NumPy}، \lr{OpenCV} و \lr{Matplotlib}) یا \lr{MATLAB}:
	
	\begin{enumerate}
		\item یک تصویر خاکستری بخوانید و هیستوگرام آن را رسم کنید.
		\item منفی، تبدیل لگاریتمی و تبدیل توانی ($\gamma=0.4$ و $\gamma=2.5$) را پیاده‌سازی کنید
		و نتیجه را کنار هم، همراه با هیستوگرام هر کدام نمایش دهید.
		\item کشش کنتراست و آستانه‌گیری (با چند مقدار $T$) را انجام دهید.
		\item $8$ بیت‌پلین تصویر را جدا و نمایش دهید و بگویید کدام‌ها اطلاعات مهم‌تری دارند.
		\item هم‌سازی هیستوگرام را \textbf{یک بار با تابع آماده} و \textbf{یک بار با پیاده‌سازی خودتان}
		انجام دهید و نتیجه‌ها را مقایسه کنید.
		\item روی یک تصویر کم‌کنتراست و یک تصویر نویزی، اثر هم‌سازی را بررسی کنید و در یک گزارش کوتاه بنویسید
		چه چیزی بهتر یا بدتر شد.
	\end{enumerate}
	
	
	% ==================================================
	\section{تمرین پیش از جلسه‌ی سوم}
	
	یک \textbf{تصویر واقعی با مشکل کیفیت} (کم‌کنتراست، تیره، روشن یا نویزی) انتخاب کنید
	و در حداکثر یک صفحه توضیح دهید:
	
	\begin{enumerate}
		\item مشکل تصویر دقیقاً چیست و هیستوگرام آن چه نشان می‌دهد؟
		\item کدام تبدیل شدتی را برای آن انتخاب می‌کنید و چرا؟
		\item تصویر و هیستوگرام قبل و بعد از اعمال تبدیل را نشان دهید.
		\item آیا مشکل به‌طور کامل حل شد؟ چه چیزی با تبدیل‌های نقطه‌ای قابل حل نبود؟
	\end{enumerate}
	
	
	% ==================================================
	\section{پیوند جلسه‌ی دوم با ادامه‌ی درس}
	
	در این جلسه تبدیل‌هایی را دیدیم که هر پیکسل را \textbf{به‌تنهایی} تغییر می‌دهند:
	
	\[
	\boxed{
		s=T(r)
	}
	\]
	
	اما دیدیم که این تبدیل‌ها برای کاهش نویز یا آشکار کردن لبه‌ها کافی نیستند.
	در جلسه‌ی سوم می‌پرسیم:
	
	\begin{center}
		\Large
		\textbf{چگونه از اطلاعات پیکسل‌های همسایه برای تعیین مقدار جدید یک پیکسل استفاده کنیم؟}
	\end{center}
	
	سپس مفاهیمی مانند:
	
	\[
	\text{همبستگی و کانولوشن}
	\]
	
	\[
	\text{فیلتر هموارساز (میانگین، گاوسی، میانه)}
	\]
	
	\[
	\text{فیلتر تیزکننده (لاپلاسین، گرادیان، \lr{Unsharp Masking})}
	\]
	
	وارد بحث خواهند شد.
	
	در واقع جلسه‌ی سوم پلی خواهد بود از:
	
	\[
	\boxed{
		\text{عملیات نقطه‌ای}
		\longrightarrow
		\text{عملیات همسایگی}
		\longrightarrow
		\text{فیلتر مکانی}
	}
	\]
	
	
	% ==================================================
	\section{خلاصه‌ی نهایی جلسه}
	
	در این جلسه دو ابزار پایه برای کار با پیکسل‌ها ساختیم.
	
	\textbf{ابزار اول:} زبان رابطه‌ی مکانی بین پیکسل‌ها:
	
	\[
	\boxed{
		N_4,\quad N_8,\quad \text{مجاورت},\quad D_e,\ D_4,\ D_8
	}
	\]
	
	\textbf{ابزار دوم:} تبدیل‌های شدتی و هیستوگرام:
	
	\[
	\boxed{
		s=T(r),
		\qquad
		h(r_k)=n_k,
		\qquad
		s_k=(L-1)\sum_{j=0}^{k}p(r_j)
	}
	\]
	
	و دیدیم که:
	
	\begin{itemize}
		\item لگاریتم و گاما تیره‌ها یا روشن‌ها را گسترش می‌دهند؛
		\item کشش کنتراست و هم‌سازی هیستوگرام، بازه‌ی شدت را بهتر استفاده می‌کنند؛
		\item هیستوگرام ابزاری ارزان برای ارزیابی تصویر است اما اطلاعات مکانی ندارد؛
		\item انتخاب تبدیل، به هدف کاربرد وابسته است؛
		\item و تبدیل‌های نقطه‌ای نمی‌توانند از اطلاعات همسایگی استفاده کنند.
	\end{itemize}
	
	برای طراحی یک روش بهبود خوب باید دائماً به سه سؤال برگردیم:
	
	\[
	\boxed{
		\begin{array}{c}
			\text{مشکل تصویر چیست؟}\\
			\text{کدام ویژگی تصویر باید تغییر کند؟}\\
			\text{چگونه نتیجه را ارزیابی کنیم؟}
		\end{array}
	}
	\]
	
	سؤال اول و دوم را امروز با کمک هیستوگرام و تبدیل‌های شدتی بررسی کردیم.
	
	سؤال مربوط به استفاده از همسایگی، موضوع اصلی جلسه‌ی سوم است.
	
	
	% ==================================================
	\section{منابع}
	
	\begin{enumerate}
		
		\item
		Rafael C. Gonzalez and Richard E. Woods,
		\textit{Digital Image Processing},
		4th Edition, Pearson, 2018
		(فصل دوم: بخش «روابط بنیادی بین پیکسل‌ها»؛
		فصل سوم: \lr{Intensity Transformations and Spatial Filtering}، بخش‌های ابتدایی).
		
		\item
		Rafael C. Gonzalez, Richard E. Woods, and Steven L. Eddins,
		\textit{Digital Image Processing Using MATLAB},
		Gatesmark Publishing.
		
		\item
		Anil K. Jain,
		\textit{Fundamentals of Digital Image Processing},
		Prentice Hall, 1989.
		
		\item
		Richard Szeliski,
		\textit{Computer Vision: Algorithms and Applications},
		Springer.
		
	\end{enumerate}
	
	
	% ==================================================
	\section*{یک جمله برای به خاطر سپردن}
	
	\begin{center}
		\Large
		\textbf{
			تبدیل شدتی یعنی هر پیکسل را به‌تنهایی تغییر دادن؛
			هیستوگرام یعنی دیدن توزیع شدت‌ها بدون دانستن مکان آنها؛
			و برای کار با لبه‌ها و نویز، باید به همسایه‌ها نگاه کرد.
		}
	\end{center}
	
\end{document}
